% ═══════════════════════════════════════════════════════════════════════
%  Stratégie 2 --- XAUUSD niveaux x10. La lecture hors échantillon n'a PAS
%  eu lieu : cette section explique pourquoi, elle ne publie aucun chiffre
%  postérieur au 2025-12-31.
% ═══════════════════════════════════════════════════════════════════════
\section{Hors échantillon}
\label{sec:x10_hors_echantillon}

\begin{warningbox}[title=\textbf{\color{white}\textsf{La tranche 2026 n'a pas été lue}}]
Le holdout de cette stratégie --- tout ce qui est postérieur au 31 décembre
2025 --- est resté verrouillé. Aucune métrique de performance n'a été calculée
dessus, par aucun des trois moteurs, et aucune ligne n'a été ajoutée au
journal de consommation du holdout, puisqu'il n'y a pas eu de lecture. Le
critère correspondant de la table de décision reste \emph{non mesuré}, donc
\textbf{en échec}.
\end{warningbox}

\lettrine[lines=2, findent=2pt, nindent=0pt]{C}{ette} section n'est pas un
aveu de travail inachevé. Ne pas lire le hors échantillon est ici une
décision, prise au vu des résultats in-sample, et c'est la bonne. Elle mérite
d'être expliquée, parce qu'elle est contre-intuitive~: on s'attend à ce qu'un
chercheur veuille toujours une mesure de plus.

\subsection{Le dispositif de gel, tel qu'il avait été prévu}

La table de décision prévoyait un protocole strict pour cette lecture~: une
section à remplir et à \emph{committer} avant que la moindre métrique
postérieure au 31 décembre 2025 soit calculée. Elle devait contenir la
configuration retenue, la fenêtre de lecture, le nombre de transactions
attendu et des intervalles de non-contradiction pour chaque métrique.

Le principe derrière ce formalisme est simple~: \textbf{une lecture hors
échantillon ne valide rien, elle peut seulement contredire}. Si l'on n'écrit
pas d'avance ce qui constituerait une contradiction, on trouvera toujours une
manière de lire le résultat qui arrange.

\begin{insightbox}
\textbf{Cette section du protocole est restée vide, et elle ne pouvait pas
être remplie.} Son premier champ est la configuration gelée. Or
\XTenPlateauPassing{} configuration ne passe le critère de plateau~: il n'y a
rien à geler. Le dispositif se bloque de lui-même, exactement comme il avait
été conçu pour le faire.
\end{insightbox}

\subsection{Trois raisons de ne pas dépenser le holdout}

\begin{description}
  \item[Le verdict est acquis in-sample.] Sur les \XTenCriteriaCount{}
  critères, \XTenCriteriaPassing{} passe. Six échecs portent une mesure
  interprétable~: espérance, profit factor, Sharpe déflaté, probabilité de
  surajustement, plateau, walk-forward, sensibilité au coût. La table de
  décision est explicite~: un seul critère en échec avec une mesure
  interprétable suffit à prononcer « ne pas déployer ». Une lecture hors
  échantillon ne peut donc rien sauver --- elle ne fait pas partie des
  conditions de réouverture.

  \item[Le holdout est un budget, pas une ressource renouvelable.] Chaque
  lecture consomme définitivement une tranche de données. La tranche 2026 de
  l'or est la seule qui existe, et elle servira à tout travail futur sur cet
  instrument --- y compris à une éventuelle stratégie issue des pistes de la
  section~\ref{sec:x10_non_deploiement}. La dépenser sur un dossier déjà
  tranché n'apprendrait rien et appauvrirait le suivant.

  \item[Le régime 2026 est hors de la distribution in-sample.] Le pas de la
  grille exprimé en \code{ATR} vaut \XTenDollarsPerAtrFirst{} au début de
  l'échantillon et \XTenDollarsPerAtrLast{} à la fin~; la tendance ne
  s'inverse pas en 2026, elle se poursuit, et le régime récent est plus fin
  que tout ce que la sélection a vu. La table de décision avait d'ailleurs
  \emph{pré-écrit} ce déclencheur~: un régime hors distribution rend la
  lecture « non concluante » par construction, avant même de la faire.
\end{description}

\subsection{Pourquoi un chiffre de 2026 serait trompeur dans les deux sens}

C'est le point qui justifie le plus directement la décision.

Supposons un résultat \textbf{positif} sur le premier trimestre 2026. Il
porterait sur quelques dizaines de transactions, dans un régime absent de
l'échantillon, face à sept années négatives et plusieurs milliers de
transactions. Le publier reviendrait à opposer un échantillon minuscule et
non représentatif à une mesure massive et cohérente --- et il se trouverait
toujours quelqu'un pour y lire un espoir.

Supposons maintenant un résultat \textbf{négatif}. Il n'ajouterait rien~: il
confirmerait ce que sept années établissent déjà, au prix d'une tranche de
données brûlée.

\begin{warningbox}
\textbf{Une mesure dont aucun résultat possible ne changerait la décision ne
doit pas être faite.} C'est le cas ici, et c'est la formulation la plus courte
de cette section. Le seul effet d'une lecture 2026 serait d'introduire, dans
un dossier par ailleurs net, un chiffre dont l'interprétation prêterait à
discussion.
\end{warningbox}

\subsection{Ce que la limite de données impose de toute façon}

Même si la lecture avait été décidée, sa portée aurait été réduite. Le
contexte dollar repose sur quatre paires de devises dont les historiques
locaux s'arrêtent au premier trimestre 2026, alors que l'or va plus loin. La
fenêtre \emph{commune aux trois moteurs} se serait donc limitée à un
trimestre. Un trimestre de barres de cinq minutes fournit beaucoup de barres,
mais peu de transactions indépendantes~: le seuil minimal de trente
transactions de la table de décision n'était même pas garanti. Dans ce cas, le
verdict prévu était « non concluant » --- et un verdict non concluant n'aurait
pas davantage autorisé le déploiement.

\subsection{Ce qui a tout de même touché des barres de 2026}

Par souci de complétude, deux mesures purement descriptives ont consommé des
barres postérieures au 31 décembre 2025, avant la campagne et sans aucun
calcul de performance~: l'\code{ATR} de cinq minutes et la fréquence de
croisement des niveaux, qui servent à établir la dérive de régime de la
section~\ref{sec:x10_grille}. Elles sont déclarées comme telles dans les
notes de recherche, et le tableau publié dans ce rapport s'arrête
volontairement à 2025 --- il ne contient aucune ligne postérieure.
